\documentclass[12pt,a4paper,reqno]{amsart}

% language
\usepackage[greek,english]{babel}
\usepackage[utf8]{inputenc}
\usepackage{alphabeta}

% change default names to greek
\addto\captionsenglish{
    \renewcommand{\contentsname}{Περιεχόμενα}
    \renewcommand{\refname}{Βιβλιογραφία}
    \renewcommand{\datename}{Ημερομηνία:}
    \renewcommand{\urladdrname}{Ιστοσελίδα}
}

% math 
\usepackage{amsmath,amsthm,amssymb,amscd}

% font
\usepackage[cal=euler]{mathalfa}
\usepackage{libertinus-type1}
% \usepackage{txfonts} % for upright greek letters
\usepackage{bm} % for bold symbols
\usepackage{bbm} % for the simply-looking bb symbols

% miscellaneous 
\usepackage{changepage} %for indenting environments
\usepackage{csquotes} % example: \textcquote{}
\usepackage{blkarray}
\setcounter{MaxMatrixCols}{20} % default for pmatrix is 10!!
\usepackage{ytableau}
\usepackage{array} %needed to increase the vertical length in a tabular

% drawing
\usepackage{tikz,tikz-cd}
\usetikzlibrary{shapes.misc, patterns, matrix, calc, intersections,positioning}
\usepackage{graphics,graphicx}
\usepackage{float} % provides enhanced control and customization options for floating objects such as figures and tables

% colors
\usepackage{xcolor}
\definecolor{darkcandyapplered}{rgb}{0.64, 0.0, 0.0}
\definecolor{midnightblue}{rgb}{0.1, 0.1, 0.44}
\definecolor{mylightblue}{HTML}{336699}
\definecolor{burntorange}{rgb}{0.8, 0.33, 0.0}
\definecolor{iceberg}{rgb}{0.44, 0.65, 0.82}
\definecolor{applegreen}{rgb}{0.55, 0.71, 0.0}
\definecolor{canaryyellow}{rgb}{1.0, 0.94, 0.0}
\definecolor{lightgray}{rgb}{0.83, 0.83, 0.83}

% hrefs
\usepackage{hyperref}
\usepackage[noabbrev,capitalize]{cleveref}
\hypersetup{
    pdftoolbar=true,        
    pdfmenubar=true,        
    pdffitwindow=false,     
    pdfstartview={FitH},  % fits the width of the page to the window
    pdftitle={},
    pdfauthor={},
    pdfsubject={},
    pdfkeywords={},
    pdfnewwindow=true,  % links in new window
    colorlinks=true,  % false: boxed links; true: colored links
    linkcolor=darkcandyapplered,   % color of internal links
    citecolor=midnightblue,  % color of links to bibliography
    urlcolor=cyan,  % color of external links
    linktocpage=true  % changes the links from the section body to the page number
    }

% geometry
\textwidth=16cm 
\textheight=21cm 
\hoffset=-55pt 
\footskip=25pt

% thm envs (you might need to change the path)
\theoremstyle{definition}
\newtheorem{theorem}{Θεώρημα}
\newtheorem{proposition}[theorem]{Πρόταση}
\newtheorem{lemma}[theorem]{Λήμμα}
\newtheorem{corollary}[theorem]{Πόρισμα}
\newtheorem{conjecture}[theorem]{Εικασία}
\newtheorem{problem}[theorem]{Πρόβλημα}
\newtheorem*{claim}{Ισχυρισμός}
\newtheorem{observation}[theorem]{Παρατήρηση}
\newtheorem{definition}[theorem]{Ορισμός}
\newtheorem{question}[theorem]{Ερώτημα}
\newtheorem*{questions}{Ερωτήματα}
\newtheorem*{que}{Ερώτημα}
\newtheorem{example}[theorem]{Παράδειγμα}
\newtheorem{exercise}{Άσκηση}

\theoremstyle{remark}
\newtheorem*{remark}{Παρατήρηση}

% fixes the correct numbering of environments
\numberwithin{theorem}{section}
\numberwithin{exercise}{section}
\numberwithin{equation}{section}

% math ops 
% fields
\newcommand{\NN}{\mathbbmss{N}} 
\newcommand{\ZZ}{\mathbbmss{Z}} 
\newcommand{\QQ}{\mathbbmss{Q}} 
\newcommand{\RR}{\mathbbmss{R}} 
\newcommand{\CC}{\mathbbmss{C}} 
\newcommand{\KK}{\mathbbmss{K}} 
\newcommand{\FF}{\mathbbmss{F}} 
\newcommand{\HH}{\mathbbmss{H}} 

% symmetric group
\newcommand{\fS}{\mathfrak{S}}  

% calligraphic 
\newcommand{\aA}{\mathcal{A}} 
\newcommand{\bB}{\mathcal{B}}
\newcommand{\cC}{\mathcal{C}}
\newcommand{\dD}{\mathcal{D}}
\newcommand{\eE}{\mathcal{E}}
\newcommand{\fF}{\mathcal{F}}
\newcommand{\hH}{\mathcal{H}}
\newcommand{\iI}{\mathcal{I}}
\newcommand{\lL}{\mathcal{L}}
\newcommand{\oO}{\mathcal{O}}
\newcommand{\pP}{\mathcal{P}}
\newcommand{\sS}{\mathcal{S}}
\newcommand{\mM}{\mathcal{M}}
\newcommand{\uU}{\mathcal{U}}

% bold
\newcommand{\bfa}{\mathbf{a}}
\newcommand{\bfe}{\mathbf{e}}
\newcommand{\bfF}{\pmb{F}}
\newcommand{\bfR}{\pmb{R}}
\newcommand{\bfv}{\mathbf{v}}
%\newcommand{\bfx}{\bm{x}}
%\newcommand{\bfx}{\mathbf{x}} 
\newcommand{\bfx}{\pmb{x}}
\newcommand{\bfX}{\pmb{X}}
\newcommand{\bfy}{\pmb{y}}
\newcommand{\bfz}{\pmb{z}}

% roman
\newcommand{\rmA}{\mathrm{A}}
\newcommand{\rmB}{\mathrm{B}}
\newcommand{\rmC}{\mathrm{C}}
\newcommand{\rmD}{\mathrm{D}} 
\newcommand{\rmH}{\mathrm{H}} 
\newcommand{\rmh}{\mathrm{h}} 
\newcommand{\rmI}{\mathrm{I}} 
\newcommand{\rmK}{\mathrm{K}}
\newcommand{\rmM}{\mathrm{M}}
\newcommand{\rmP}{\mathrm{P}}  
\newcommand{\rmp}{\mathrm{p}}  
\newcommand{\rmQ}{\mathrm{Q}}  
\newcommand{\rmR}{\mathrm{R}}
\newcommand{\rmS}{\mathrm{S}}
\newcommand{\rmT}{\mathrm{T}}
\newcommand{\rmU}{\mathrm{U}}
\newcommand{\rmV}{\mathrm{V}}
\newcommand{\rmY}{\mathrm{Y}}
\newcommand{\rmZ}{\mathrm{Z}}
\newcommand{\rmz}{\mathrm{z}}

% arrows and symbols 
\renewcommand{\to}{\rightarrow}
\newcommand{\toto}{\longrightarrow}
\newcommand{\mapstoto}{\longmapsto}
\newcommand{\then}{\Rightarrow}
\newcommand{\IFF}{\Leftrightarrow}
\newcommand{\tl}{\tilde}
\newcommand{\wtl}{\widetilde}
\newcommand{\ol}{\overline}
\newcommand{\ul}{\underline}
\newcommand{\oldemptyset}{\emptyset}
\renewcommand{\emptyset}{\varnothing}
\DeclareMathSymbol{\Arg}{\mathbin}{AMSa}{"39} % for arguments 
\newcommand{\onto}{\ensuremath{\twoheadrightarrow}}
\newcommand{\tle}{\trianglelefteq}
\newcommand{\tge}{\trianglerighteq}

% custom ops
\newcommand{\rmKer}{\mathrm{Ker}}
\newcommand{\rmIm}{\mathrm{Im}}
\newcommand{\lcm}{\mathrm{lcm}}
\newcommand{\GL}{\mathrm{GL}}
\newcommand{\ev}{\mathrm{ev}}

% absolute value symbol
\usepackage{mathtools} 
\DeclarePairedDelimiter\abs{\lvert}{\rvert}%
\DeclarePairedDelimiter\norm{\lVert}{\rVert}%
\makeatletter
\let\oldabs\abs
\def\abs{\@ifstar{\oldabs}{\oldabs*}}

% tensor symbol
\newcommand{\tensor}[1]{%
  \mathbin{\mathop{\otimes}\limits_{#1}}%
}

% permutation cycle notation
\ExplSyntaxOn
\NewDocumentCommand{\cycle}{ O{\;} m }
 {
  (
  \alec_cycle:nn { #1 } { #2 }
  )
 }

\seq_new:N \l_alec_cycle_seq
\cs_new_protected:Npn \alec_cycle:nn #1 #2
 {
  \seq_set_split:Nnn \l_alec_cycle_seq { , } { #2 }
  \seq_use:Nn \l_alec_cycle_seq { #1 }
 }
\ExplSyntaxOff

% setminus symbol
\newcommand{\mysetminusD}{\hbox{\tikz{\draw[line width=0.6pt,line cap=round] (3pt,0) -- (0,6pt);}}}
\newcommand{\mysetminusT}{\mysetminusD}
\newcommand{\mysetminusS}{\hbox{\tikz{\draw[line width=0.45pt,line cap=round] (2pt,0) -- (0,4pt);}}}
\newcommand{\mysetminusSS}{\hbox{\tikz{\draw[line width=0.4pt,line cap=round] (1.5pt,0) -- (0,3pt);}}}
\newcommand{\sm}{\mathbin{\mathchoice{\mysetminusD}{\mysetminusT}{\mysetminusS}{\mysetminusSS}}}

% miscellaneous commands
\newcommand{\defn}[1]{{\color{mylightblue}{#1}}}
\newcommand{\toDo}{{\bf\color{red} TODO}}
\newcommand{\toCite}{{\bf\color{green} CITE}}
\newcommand*{\vertbar}{\rule[-1ex]{0.5pt}{2.5ex}} % for matrices with column vectors
\newcommand*{\horzbar}{\rule[.5ex]{2.5ex}{0.5pt}} % for matrices with row vectors
\newcommand{\myblue}[1]{{\color{iceberg}{#1}}}
\newcommand{\myorange}[1]{{\color{burntorange}{#1}}}
\newcommand{\mygreen}[1]{{\color{applegreen}{#1}}}
\newcommand{\myred}[1]{{\color{darkcandyapplered}{#1}}}

% ferrer's diagram
\newcommand{\fdiagram}[1]{
    \begin{tikzpicture}[scale=.7]
        \fill foreach \Z [count=\Y] in {#1}
        {foreach \X in {1,...,\Z} 
        {(\X,-\Y) circle[radius=3pt]}};
    \end{tikzpicture}
}

%
\makeatletter
\def\mathcolor#1#{\@mathcolor{#1}}
\def\@mathcolor#1#2#3{%
  \protect\leavevmode
  \begingroup
    \color#1{#2}#3%
  \endgroup
}
\makeatother

% 
\newenvironment{nouppercase}{%
  \let\uppercase\relax%
  \renewcommand{\uppercasenonmath}[1]{}}{}

% titlepage
\title{226: Θεωρία Δακτυλίων και Modules}
\author[Β.~Δ. Μουστακας]{Βασίλης Διονύσης Μουστάκας \\ Πανεπιστήμιο Κρήτης}
\date{26 Φεβρουαρίου 2026}
% \urladdr{\href{https://sites.google.com/view/vasmous}{https://sites.google.com/view/vasmous}}

\begin{document}

\begingroup
\def\uppercasenonmath#1{} % this disables uppercase title
\let\MakeUppercase\relax % this disables uppercase authors
\maketitle
\endgroup

\setcounter{section}{3}
% \setcounter{theorem}{11}
% \setcounter{equation}{}
\begin{center}
    \textbf{3. Πολυωνυμικοί δακτύλιοι και παραγοντοποίηση
} 
\end{center}

Έστω $R$ ένας μεταθετικός δακτύλιος. Ένας \defn{πολυωνυμικός δακτύλιος σε μία μεταβλητή} είναι ο ένας μεταθετικός δακτύλιος $R[x]$ που προκύπτει από τον $R$ \textquote{προσθέτοντας} ένα νέο στοιχείο $x$ το οποίο μετατίθεται με κάθε στοιχείο του $R$, δηλαδή
\[
xr = rx
\]
για κάθε $r \in R$ και είναι \textquote{ελέυθερο σχέσεων} στο $R$, δηλαδή αν 
\[
a_0 + a_1x + \cdots + a_nx^n = 0
\]
για κάποια $a_0, a_1, \dots, a_n \in R$ είναι μία \textquote{σχέση} στο $R$, τότε 
\[
a_0 = a_1 = \cdots = a_n = 0.
\]
Για την απόδειξη της ύπαρξης ενός τέτοιου δακτύλιου παραπέμπουμε στο \cite[Θεώρημα~2.2.1]{VDEMT13}.

Τα στοιχεία αυτού του δακτύλιου, τα οποία ονομάζονται \defn{πολυώνυμα} είναι 
\begin{itemize}
    \item τα στοιχεία του $R$ (\defn{σταθερά πολυώνυμα}),
    \item δυνάμεις του $x$ (\defn{μονώνυμα}), και
    \item πεπερασμένοι \textquote{$R$-γραμμικοί συνδυασμοί} δυνάμεων του $x$, δηλαδή 
    \[
    a, \ a + bx, \ a + bx + cx^2, \ \dots 
    \]
    για $a, b$ και $c \in R$.
\end{itemize}
Πιο συγκεκριμένα, 
\[
R[x] = \{ a_0 + a_1x + \cdots + a_nx^n : n \in \NN, a_0, a_1, \dots, a_n \in R\}.
\]
Αν 
\[
f = f(x) = a_0 + a_1x + \cdots + a_nx^n \ \in \ R[x]
\]
με $a_n \neq 0$, τότε το 
\begin{itemize}
    \item $a_0$ ονομάζεται \defn{σταθερός όρος} του $f$
    \item $a_n$ ονομάζεται \defn{μεγιστοβάθμιος συντελεστής} του $f(x)$
    \item $n$ ονομάζεται \defn{βαθμός} του $f$ και συμβολίζεται με $\deg(f(x))$.
\end{itemize}

Πολλές από τις αριθμητικές ιδιότητες του δακτύλιου $R$ που συζητήσαμε στις προηγούμε\-νες δύο ενότητες μεταφέρονται και στον πολυωνυμικό δακτύλιο $R[x]$. Για παράδειγμα, αν ο $R$ είναι ακέραια περιοχή, τότε το ίδιο ισχύει και για τον $R[x]$. Πράγματι, έστω $f, g \in R[x]$ δύο μη σταθερά πολυώνυμα βαθμού $n$ και $m$, αντίστοιχα. Τότε 
\begin{align*}
    fg 
    &= (a_0 + a_1x + \cdots + a_nx^n)(b_0 + b_1x + \cdots + b_mx^m) \\ 
    &= a_nb_m x^{n+m} + \ \left(\text{μονώνυμα βαθμού $< n + m$}\right) \neq 0, 
\end{align*}
καθώς $a_nb_m \neq 0$, διότι ο $R$ είναι ακέραια περιοχή. 

Ο παραπάνω υπολογισμός μας πληροφορεί ότι αν ο $R$ είναι μια ακέραια περιοχή, τότε 
\begin{equation}
    \label{eq:degree}
    \deg(fg) = \deg(f) + \deg(g)
\end{equation}
για κάθε $f, g \in R[x]$. Γενικότερα, ισχύει η ανισότητα 
\[
\deg(fg) \leq \deg(f) + \deg(g)
\]
(γιατί;). Για παράδειγμα, αν $R = \ZZ_6$, τότε 
\[
(2)\cdot(3x) = 6x = 0
\]
στο $\ZZ_6[x]$.

Μια ακόμη πληροφορία που παίρνουμε από τον υπολογισμό αυτόν είναι ότι αν ο $R$ είναι ακέραια περιοχή, τότε 
\[
\left(R[x]\right)^\times = R^\times.
\]
Πράγματι, έστω $f \in \left(R[x]\right)^\times$. Τότε $fg = 1$, για κάποιο  $g \in R[x]$. Από την Ταυτότητα \eqref{eq:degree} έπεται ότι 
\[
0 = \deg(1) = \deg(fg) = \deg(f) + \deg(g).
\]
Συνεπώς, $f, g \in R$ και γι αυτό $f \in R^\times$. Άρα, $\left(R[x]\right)^\times \subseteq R^\times$ και η άλλη κατεύθυνση είναι προφανής (γιατί;).

Το στοιχείο $x$ μπορούμε να το φανταστούμε ως μεταβλητή με την εξής έννοια.
\begin{proposition}
    \label{prop:evaluation_hom}
    Έστω $S$ ένας μεταθετικός δακτύλιος που περιέχει τον $R$. Για κάθε $s \in S$, υπάρχει μοναδικός ομομορφισμός δακτυλίων
    \begin{align*}
        \ev_s : R[x] &\to S \\
                   x &\mapsto s
    \end{align*}
    ο οποίος ονομάζεται \defn{ομομορφισμός εκτίμησης} (evaluation homomorphism).
\end{proposition}

Για παράδειγμα, για $R = S = \RR$ και $f = x^2 + 1 \in \RR[x]$, η εκτίμηση του $f$ στο $1$ είναι 
\[
\ev_1(f) = \ev_1(x^2 + 1) = \ev_1(x)^2 + \ev_1(1) = 1 + 1 = 2,
\]
σε συμφωνία με τον πιο οικείο συμβολισμό $f(1)$.

Τι γίνεται αν θέλουμε να \textquote{εκτιμήσουμε} το $x$ σε ένα στοιχείο από κάποιον δακτύλιο που δε σχετίζεται με τον $R$; Έστω $S$ ένας μεταθετικός δακτύλιος και $\phi : R \to S$ ένας ομομορφισμός δακτυλίων. Για κάθε $s \in S$, υπάρχει μοναδικός ομομορφισμός δακτυλίων $\tilde{\phi}_s : R[x] \to S$ με $\tilde{\phi}_s(x) = s$, τέτοιος ώστε το διάγραμμα 
\[
\begin{tikzcd}
& R \arrow[dl, "\iota"'] \arrow[dr, "\phi"] & \\
R[x] \arrow[rr, "\tilde{\phi}_s"'] & & S
\end{tikzcd}
\]
να μετατίθεται, δηλαδή $\phi = \tilde{\phi}_s \circ \iota$, όπου $\iota : R \to R[x]$ είναι ο \textquote{φυσικός} μονομορφισμός που απεικονίζει τα στοιχεία του $R$ στα σταθερά πολυώνυμα του $R[x]$. Στην περίπτωση αυτή, έχουμε 
\[
\tilde{\phi}_s(a_0 + a_1x + \cdots + a_nx^n) = 
\phi(a_0) + \phi(a_1)s + \cdots + \phi(a_n)s^n
\]
για κάθε $a_0 + a_1x + \cdots + a_nx^n \in R[x]$.

Για παράδειγμα, έστω $R = \RR, S = \rmM_2(\RR)$ και θεωρούμε τον ομομορφισμό δακτυλίων 
\begin{align*}
        \phi : \RR &\to \rmM_2(\RR) \\
                1  &\mapsto \begin{pmatrix}
                    1 & 0 \\ 0 & 1
                \end{pmatrix}.
\end{align*}
Τότε, για $s = \left(\begin{smallmatrix}
0 & 1 \\ 0 & 0
\end{smallmatrix}\right)$ υπολογίζουμε 
\begin{align*}
\tilde{\phi}_{\left(\begin{smallmatrix}
0 & 1 \\ 0 & 0
\end{smallmatrix}\right)}(-1 + 2x + x^3) &= 
- \phi(1) + \phi(2)\begin{pmatrix}
    0 & 1 \\ 0 & 0
\end{pmatrix} + \begin{pmatrix}
    0 & 1 \\ 0 & 0
\end{pmatrix}^3 \\
&= 
-\begin{pmatrix}
    1 & 0 \\ 0 & 1
\end{pmatrix}
+ 2\begin{pmatrix}
    1 & 0 \\ 0 & 1
\end{pmatrix}\begin{pmatrix}
    0 & 1 \\ 0 & 0
\end{pmatrix} \\ 
&= 
\begin{pmatrix}
    -1 & 0 \\ 0 & -1
\end{pmatrix}
+ \begin{pmatrix}
    0 & 2 \\ 0 & 0
\end{pmatrix} \\ 
&= \begin{pmatrix}
    -1 & 2 \\
    0  & -1
\end{pmatrix}.
\end{align*}

\begin{remark}
    Η Πρόταση \ref{prop:evaluation_hom} είναι ειδική περίπτωση της γενικότερης κατασκευής που μόλις περιγράψαμε, όπου η $\phi$ είναι απλώς η ταυτοτική απεικόνιση. Αυτό είναι ένα παράδειγμα \emph{καθολικής ιδιότητας}, η οποία μας πληροφορεί ότι ο πολυωνυμικός δακτύλιος σε μία μεταβλητή είναι ο \textquote{μοναδικός} μεταθετικός δακτύλιος ο οποίος διαθέτει μια \textquote{μεταβλητή} την οποία μπορούμε να εκτιμήσουμε. Μας πληροφορεί, δηλαδή, αυτό που \emph{κάνει} ένας πολυωνυμικός δακτύλιος και όχι αυτό που \emph{είναι}.
\end{remark}

Έχοντας ορίσει τον πολυωνυμικό δκατύλιο σε μία μεταβλητή, μπορούμε να ορίσουμε τον πολυωνυμικό δακτύλιο $R[x_1, x_2, \dots, x_n]$ πολλών μεταβλητών, επαγωγικά ως εξής 
\[
R[x_1, x_2, \dots, x_n] \coloneqq \left(R[x_1, \dots, x_{n-1}]\right)[x_n].
\]
Τα στοιχεία αυτού του δακτύλιου είναι 
\begin{itemize}
    \item τα στοιχεία του $R$ (\defn{σταθερά πολυώνυμα}),
    \item γινόμενα δυνάμεων των $x_1, x_2, \dots, x_n$, δηλαδή
    \[
    x_1^{\alpha_1}x_2^{\alpha_2}\cdots x_n^{\alpha_n},
    \]
    που ονομάζονται \defn{μονώνυμα} με \defn{εκθέτες} $\alpha_1, \alpha_2, \dots, \alpha_n \in \NN$, και
    \item πεπερασμένοι \textquote{$R$-γραμμικοί συνδυασμοί} μονωνύμων.
\end{itemize}

Πιο συγκεκριμένα, 
\[
R[x_1, x_2, \dots, x_n] = 
\left\{
    \sum_{\alpha_1, \alpha_2, \dots, \alpha_n \in \NN} c_{\alpha_1, \alpha_2, \dots, \alpha_n} x_1^{\alpha_1}x_2^{\alpha_2}\cdots x_n^{\alpha_n} : n \in \NN, c_{\alpha_1, \alpha_2, \dots, \alpha_n} \in R
\right\}.
\]
Το άθροισμα των εκθετών ενός μονωνύμου ονομάζεται \defn{βαθμός} του και ο βαθμός ενός πολυω\-νύμου πολλών μεταβλητών ορίζεται ως ο μέγιστος βαθμός ενός μονωνύμου στο ανάπτυγμά του. Για παράδειγμα, το 
\[
4 - 3x_1x_2 + 2x_2^2 - x_1^2x_2^3 \ \in \ \RR[x_1,x_2]
\]
είναι ένα πολυώνυμο δύο μεταβλητών βαθμού 5, ενώ ως πολυώνυμο στην μεταβλητή $x_2$ με συντελεστές στο $\RR[x_1]$ έχει την μορφή 
\[
\mathcolor{gray}{4} - \mathcolor{gray}{3x_1}x_2 + \mathcolor{gray}{2}x_2^2 - \mathcolor{gray}{x_1^2}x_2^3
\]
και βαθμό 3.

Όμοια με την περίπτωση του πολυωνυμικού δακτύλιου μίας μεταβλητής\footnote{Ο πολυωνυμικός δακτύλιος πολλών μεταβλητών ικανοποιεί και αυτός μια καθολική ιδιότητα όμοια με την περίπτωση μίας μεταβλητής (ποιά είναι η διατύπωσή της;).}, έτσι και εδώ μπορούμε να φανταστούμε τα $x_1, x_2, \dots, x_n$ ως μεταβλητές με την εξής έννοια.
\begin{proposition}
    \label{prop:evaluation_hom_multivariable}
    Έστω $S$ ένας μεταθετικός δακτύλιος που περιέχει τον $R$. Για κάθε $s_1, s_2, \dots,$ $s_n \in S$, υπάρχει μοναδικός ομομορφισμός δακτυλίων 
    \begin{align*}
        \ev_{s_1, s_2, \dots, s_n} : R[x_1, x_2, \dots, x_n] &\to S \\
                   x_i &\mapsto s_i
    \end{align*}
    για κάθε $1 \le i \le n$.
\end{proposition}

Ας στρέψουμε την προσοχή μας στην αριθμητική των πολυωνυμικών δακτυλίων.
\begin{theorem}
    \label{thm:polynomial_rings_are_ED}
    Έστω $F$ σώμα. Αν $f, g \in F[x]$ με $g \neq 0$, τότε υπάρχουν (μοναδικά) $q, r \in F[x]$ τέτοια ώστε 
    \[
    f = qg + r
    \]
    όπου $r=0$ ή $\deg(r) < \deg(g)$. Ειδικότερα, ο $F[x]$ είναι Ευκλείδια περιοχή και κατά συνέπεια περιοχή κύριων ιδεωδών και περιοχή μονοσήμαντης παραγοντοποίησης.
\end{theorem}

Για μία απόδειξη παραπέμπουμε στο \cite[Θεώρημα~2.3.3]{VDEMT13}. Το Θεώρημα \ref{thm:polynomial_rings_are_ED} ισχύει και με την ασθενέστερη υπόθεση ότι το $F$ είναι ακέραια περιοχή αντί για σώμα και ότι ο μεγιστοβάθμιος συντελεστής του $g$ είναι αντιστρέψιμο στοιχείο, αντί για $g \neq 0$ (βλ. \cite[Θεώρημα~2.3.5]{VDEMT13}). 

\begin{remark} 
    Παραδόξως, ισχύει και το εξής αντίστροφο του Θεωρήματος \ref{thm:polynomial_rings_are_ED}: Αν ο $R[x]$ είναι περιοχή κύριων ιδεωδών, τότε ο $R$ πρέπει να είναι σώμα. Πράγματι, αν ο $R[x]$ είναι περιοχή κύριων ιδεωδών, τότε ο $R$ πρέπει να είναι ακέραια περιοχή (ως υποδακτύλιος ακέραιας περιοχής). Όμως, $R[x]/(x) \cong R$ και γι αυτό από την Πρόταση 1.11 έπεται ότι το $(x)$ είναι πρώτο ιδεώδες. Από το Λήμμα 2.14 αυτό σημαίνει ότι είναι και μεγιστικό και γι αυτό ο $R \cong R[x]/(x)$ είναι σώμα από την Πρόταση 1.10.
\end{remark}

Από την παρατήρηση έπεται ότι οι πολυωνυμικοί δακτύλιοι 
\[
\ZZ[x], \ F[x_1, x_2], \ \dots, \ F[x_1, x_2, \dots, x_n]
\]
δεν είναι περιοχές κύριων ιδεωδών, διότι ο δακτύλιος των ακεραίων και ο πολυωνυμικός δακτύλιος μιας μεταβλητής δεν είναι σώματα. 

\begin{definition}
    \label{def:ideals_generated_by}
    Έστω $R$ ένας μεταθετικός δακτύλιος και $a_1, a_2, \dots, a_n \in R$. Το ιδεώδες 
    \[
    (a_1, a_2, \dots, a_n) \coloneqq (a_1) + (a_2) + \cdots + (a_n)
    \]
    ονομάζεται το \defn{ιδεώδες που παράγεται από} τα στοιχεία $a_1, a_2, \dots, a_n$. Κάθε ιδεώδες του $R$ αυτής της μορφής ονομάζεται \defn{πεπερασμένα παραγόμενο} (finitely generated).
\end{definition}

Το ιδεώδες που παράγεται από τα $a_1, a_2, \dots, a_n$ είναι το μικρότερο ιδεώδες του $R$ που περιέχει τα $a_1, a_2, \dots, a_n$, δηλαδή 
\[
(a_1, a_2, \dots, a_n) = \bigcap_{\substack{\text{$I$ ιδεώδες του $R$} \\ a_1, a_2, \dots a_n \ \in \ I}} I
\]
(βλ. \cite[Πρόταση~2.5.9]{VDEMT13}). 

\begin{example}
    \label{ex:Z[x]_is_not_a_PID}
    Ας δούμε γιατί το $\ZZ[x]$ δεν είναι περιοχή κύριων ιδεωδών στην πράξη, δείχνοντας ότι το ιδεώδες $(2, x)$ που παράγεται από το $2$ και το $x$ δεν είναι κύριο. Έστω $f \in \ZZ[x]$ τέτοιο ώστε 
    \[
    (f) = (2, x) = \{2p(x) + xq(x) : p, q \in \ZZ[x]\}.
    \]
    Επειδή $2 \in (f)$, έπεται ότι 
    \[
    2 = f(x)p(x),
    \]
    για κάποιο $p \in \ZZ[x]$. Από την Ταυτότητα \eqref{eq:degree} έπεται ότι 
    \[
    \deg(f) + \deg(p) = 0,
    \]
    και γι αυτό $f, p \in \ZZ$ με τις εξής επιλογές
    \[
    \renewcommand{\arraystretch}{1.3} 
    \begin{tabular}{c|c|c|c|c}
    $f$ & $1$ & $2$ & $-1$ & $-2$\\ \hline
    $p$ & $2$ & $1$ & $-2$ & $-1$
    \end{tabular}
    \]
    Αν $f = \pm 1$, τότε $(f) = (2,x) = \ZZ[x]$, το οποίο είναι αδύνατο, διότι 
    \[
    3 \notin (2, x).
    \]
    Αν $f = \pm 2$, τότε επειδή $x \in (f) = (2)$ έπεται ότι 
    \[
    x = 2q(x)
    \]
    για κάποιο $q \in \ZZ[x]$, το οποίο είναι και αυτό αδύνατο. 
\end{example}


\begin{thebibliography}{99}
    \bibitem{VDEMT13}
    Δ. Α. Βάρσος, Δ. Ι. Δεριζιώτης, Ι. Π. Εμμανουήλ, Μ. Π. Μαλιάκας και Ο. Π. Ταλέλλη, 
    \emph{Μια Εισαγωγή στην Άλγεβρα}, 
    τρίτη έκδοση, Εκδόσεις Σοφία, 2013.
\end{thebibliography}
\end{document}